\documentclass[../CLT_SYK.tex]{subfiles}
\setcounter{section}{1}

\begin{document}

\section{\textcolor{red}{Polynomial chaos expansion}\ \textcolor{blue}{Gaussian-chaos decomposition}}
\label{sec:chaos-decomposition}
We call $\alpha=(\alpha_I)_{I\in\mathcal I_N} \in (\mathbb N\cup\{0\})^{\mathcal I_N}$  a multi-index. We define the order of $\alpha$ as $|\alpha| := \sum_{I\in\mathcal I_N} \alpha_I $.
 \textcolor{blue}{We use the probabilists' Hermite polynomials}
 \textcolor{red}{For $k\in\mathbb Z_{\geq0}$, the probabilists' Hermite
 polynomial of degree $k$ is defined by
 \[
   {\rm He}_k(x):=(-1)^k e^{x^2/2}\frac{d^k}{dx^k}e^{-x^2/2}.
 \]
 These polynomials are orthogonal with respect to the standard Gaussian
 measure; in particular, ${\rm He}_0(x)=1$, ${\rm He}_1(x)=x$, and
 ${\rm He}_2(x)=x^2-1$.  Their multivariate products satisfy}
\begin{equation*}
  H_\alpha(G):=\prod_{I\in\calI_N}{\rm He}_{\alpha_I}(g_I),
  \qquad
  \E[H_\alpha H_\eta]
  =\alpha!\,\mathbf 1_{\{\alpha=\eta\}},
\end{equation*}
where $\alpha!:=\prod_I\alpha_I!$\textcolor{red}{, and
$\mathbf 1_{\{\alpha=\eta\}}$ is the indicator function}. Since the $H_\alpha$ form a complete
orthogonal basis of $L^2(\gamma_{\mathcal I_N})$, where
$\gamma_{\mathcal I_N}$ denotes \textcolor{red}{the} standard Gaussian measure on
$\mathbb R^{\mathcal I_N}$, we have
\begin{align}\label{eq: chaos}
Z_N(\beta;G) =\textcolor{red}{\sum_{k=0}^\infty \sum_{\alpha:|\alpha|=k} \frac{\E[Z_N(\beta;G)H_\alpha(G)]}{\alpha!} H_\alpha(G)}\quad \textcolor{blue}{\sum_{k=0}^\infty \sum_{\alpha:|\alpha|=k} \frac{\E[Z_N(\beta;G')H_\alpha(G')]}{\alpha!} H_\alpha(G)},
\end{align}
where the second sum runs over all multi-indices $\alpha$ such that $\sum_{I} \alpha_I = k$\textcolor{blue}{, and the expectation is taken over an independent copy $G'$ of $G$}.  Let $P_r$ denote the orthogonal projection in
$L^2(\gamma_{\mathcal I_N})$ onto the $r$-th \textcolor{blue}{homogeneous Gaussian} chaos \textcolor{red}{component}.
Then, we can write
\[
  P_r Z_N(\beta;G)
  =\textcolor{red}{\sum_{|\alpha|=r}
  	\frac{\E[Z_N(\beta;G)H_\alpha(G)]}{\alpha!}H_\alpha(G)}\
  	\textcolor{blue}{\sum_{|\alpha|=r}
  \frac{\E[Z_N(\beta;G')H_\alpha(G')]}{\alpha!}H_\alpha(G)}.
\]
\textcolor{red}{The aim of this section is to prove the following proposition for the normalized partition function.
\begin{proposition}\label{prop:exact-second-chaos}
	For every $\beta\in \mathbb R$,
	\[
	P_1Y_N(\beta)=0\quad \text{and} \quad P_2 Y_N(\beta)
	=A_N(\beta)Q_N,
	\]
	where $A_N(\beta)$ and $Q_N$ are defined by \eqref{eq:finite-N-FTW-coefficient}.
\end{proposition}
}
\textcolor{blue}{
Since $\E Z_N(\beta)$ is deterministic, we also have
\[
  P_rY_N(\beta)=\frac{P_rZ_N(\beta)}{\E[Z_N(\beta)]}.
\]
Recall $A_N(\beta):=\frac{\beta}{2}\,\partial_\beta\log \E Z_N(\beta)$ from \eqref{eq:finite-N-FTW-coefficient}.
 The next lemma proves the exact radial decomposition
\begin{equation*}
  Y_N(\beta)-1=A_N(\beta)Q_N+\mathcal R_N(\beta),
  \qquad
  \mathcal R_N:=\sum_{r\geq3}P_r Y_N.
\end{equation*}}


\textcolor{blue}{\subsection{Gauge symmetry and the parity selection rule}}
 \textcolor{red}{We call} \textcolor{blue}{
Call} a multi-index $\alpha$ \emph{even} if \textcolor{red}{$\sum_{I\colon x \in I}
\alpha_I$ is even for every $x\in[N]$, and odd otherwise}\ \textcolor{blue}{for every $x\in[N],$
\begin{equation}\label{eq:even-multi-index}
  \sum_{I\colon x \in I}
  \alpha_I\equiv0\pmod2.
\end{equation}
Otherwise, we call $\alpha$ odd}. \textcolor{red}{We write $\mathfrak{S}_N$ for the symmetric group of degree $N$.}
\begin{lemma}\label{lem: parity rule lemma}
    For $s=(s_1,\ldots,s_N)\in\{\pm1\}^N$, define
$\chi_s(I):=\prod_{x\in I}s_x$.
\textcolor{red}{For $\pi\in\mathfrak S_N$ and $I=\{i_1,\ldots,i_q\}\in\calI_N$ with $i_1<\cdots<i_q$, we also
define $\epsilon_\pi(I)\in\{\pm1\}$ by
\[
  \mathrm i^{q/2}\psi_{\pi(i_1)}\cdots\psi_{\pi(i_q)}
  =\epsilon_\pi(I)\psi_{\pi(I)},
\]
whose existence follows from the anticommutation relation \eqref{eq: clifford relation}.
Set
\[
  T_{\pi,s}G
  :=\bigl(\chi_s(I)\epsilon_\pi(I)g_{\pi(I)}\bigr)_{I\in\calI_N}.
\]
Then, for every $\beta\in\mathbb R$ and $G\in\mathbb R^{\calI_N}$,
\begin{equation}\label{eq:signed-permutation-invariance}
  Z_N(\beta;T_{\pi,s}G)=Z_N(\beta;G).
\end{equation}
}
\textcolor{blue}{We have
\begin{equation}\label{eq:sign-gauge-invariance}
  Z_N(\beta;G)=Z_N\bigl(\beta;(\chi_s(I)g_I)_I\bigr).
\end{equation}}
    In particular, if $\alpha$ is odd, then
   \begin{equation}\label{eq:odd-coefficient-zero}
  \E[Z_N(\beta;G)H_\alpha(G)]=0.
   \end{equation}
    
\end{lemma}
\begin{proof}
We define \textcolor{blue}{the Clifford volume element:} 
\[
  \Gamma
  := 
  \psi_1\cdots\psi_N.
\]
Note that $\Gamma^*\Gamma$ is \textcolor{red}{the} identity \textcolor{red}{matrix} \textcolor{blue}{map}, so $\Gamma$ is unitary. Since $N$ is even, moving $\psi_j$ through the other $N-1$ Majorana
\textcolor{red}{matrices} \textcolor{blue}{generators} produces the sign $(-1)^{N-1}=-1$.  Hence, for any $j\in[N],$
\[
  \Gamma\psi_j\Gamma^*=  -\psi_j \Gamma \Gamma^* =-\psi_j.
\]
\textcolor{red}{For a real unit vector $v=(v_1,\ldots,v_N)$, set
\[
  C_v:=\sum_{j=1}^N v_j\psi_j,
  \qquad U_v:=\Gamma C_v.
\]
The anticommutation relation~\eqref{eq: clifford relation} gives
\[
  C_v^*=C_v,\qquad C_v^2=\mathbb I,\qquad
  C_v\psi_j+\psi_jC_v=2v_j\mathbb I\quad(j\in[N]).
\]
Thus $C_v$ and $U_v$ are unitary, and
\begin{equation}\label{eq:reflection-conjugation}
  U_v\psi_jU_v^*
  =\Gamma(2v_jC_v-\psi_j)\Gamma^*
  =\psi_j-2v_jC_v.
\end{equation}
In particular, taking $v=e_x$, where $e_x$ is the $x$-th
standard basis vector of $\mathbb R^N$, gives
$U_x:=\Gamma\psi_x$ and}
\textcolor{blue}{For each $x\in[N]$,  consider a single-coordinate reflection
$$U_x:=\Gamma\psi_x.$$
This is unitary because $U_x^* U_x=\ \psi_x^* \Gamma^*\Gamma\psi_x  =\mathbb I$.  The  Clifford relations \eqref{eq: clifford relation} give
\[
  \psi_x\psi_j\psi_x
  =
  \begin{cases}
    \psi_x,&j=x,\\
    -\psi_j,&j\neq x.
  \end{cases}
\]
Hence, we obtain
}
\begin{align}\label{eq: conjugation rule}
      U_x\psi_jU_x^*
  =
  \begin{cases}
    -\psi_x,&j=x,\\
    \psi_j,&j\neq x.
  \end{cases}
\end{align}



Let $S_s:=\{x\in[N]:s_x=-1\}$ and let $U_s$ be the product of the unitaries
$U_x$ over $x\in S_s$:
$$U_s:= \prod_{x\in S_s} U_x,$$
where the product is in any fixed order.  Applying \eqref{eq: conjugation rule} successively yields that for any $j\in[N],$
\[
  U_s\psi_jU_s^*=s_j\psi_j.
\]
\textcolor{red}{For distinct $a,b\in[N]$, taking
$v=(e_a-e_b)/\sqrt2$ in~\eqref{eq:reflection-conjugation} gives
\begin{equation}\label{eq:unitary-transposition}
 V_{ab}:=\Gamma\frac{\psi_a-\psi_b}{\sqrt2},
\qquad
V_{ab}\psi_jV_{ab}^*
=\begin{cases}
	\psi_b,&j=a,\\
	\psi_a,&j=b,\\
	\psi_j,&j\notin\{a,b\}.
\end{cases}
\end{equation}
Writing $\pi=\tau_m\circ\cdots\circ\tau_1$ as a product of
transpositions and setting $V_\pi:=V_{\tau_m}\cdots V_{\tau_1}$,
we obtain $V_\pi\psi_jV_\pi^*=\psi_{\pi(j)}$ by \eqref{eq:unitary-transposition}.
For $\pi=\mathrm{id}$, we take $V_\pi=\mathbb I$.
The unitary $V_{\pi,s}:=V_\pi U_s$ therefore satisfies
\[
  V_{\pi,s}\psi_jV_{\pi,s}^*=s_j\psi_{\pi(j)}.
\]
Recall that $\psi_I=\mathrm i^{q/2}\psi_{i_1}\cdots\psi_{i_q}$
for $I=\{i_1,\ldots,i_q\}$ with $i_1<\cdots<i_q$.
By the definitions of $\chi_s(I)$ and $\epsilon_\pi(I)$,
\begin{align*}
  V_{\pi,s}\psi_IV_{\pi,s}^*
  &=\mathrm i^{q/2}
    (V_{\pi,s}\psi_{i_1}V_{\pi,s}^*)\cdots
    (V_{\pi,s}\psi_{i_q}V_{\pi,s}^*)\\
  &=\chi_s(I)\epsilon_\pi(I)\psi_{\pi(I)}.
\end{align*}
Since $I\mapsto\pi(I)$ is a permutation of $\calI_N$, we have
\begin{align*}
  V_{\pi,s}H_N(T_{\pi,s}G)V_{\pi,s}^*
  &=\sigma_N\sum_{I\in\calI_N}
    \chi_s(I)\epsilon_\pi(I)g_{\pi(I)}
    V_{\pi,s}\psi_IV_{\pi,s}^*\\
  &=\sigma_N\sum_{I\in\calI_N}
    \bigl(\chi_s(I)\epsilon_\pi(I)\bigr)^2
    g_{\pi(I)}\psi_{\pi(I)}
    =H_N(G).
\end{align*}
Hence, we conclude that
\[
Z_N(\beta;T_{\pi,s}G)
=\tr\!\left(
V_{\pi,s}e^{-\beta H_N(T_{\pi,s}G)}V_{\pi,s}^*
\right)
=\tr e^{-\beta H_N(G)}
=Z_N(\beta;G),
\]
because the normalized trace is invariant under unitary conjugation
and
\[
V_{\pi,s}e^{-\beta H_N(T_{\pi,s}G)}V_{\pi,s}^*
=e^{-\beta V_{\pi,s}H_N(T_{\pi,s}G)V_{\pi,s}^*}
=e^{-\beta H_N(G)}.
\]
This proves~\eqref{eq:signed-permutation-invariance}.}

\textcolor{blue}{Recall that $\psi_I=\mathrm i^{q/2}\psi_{i_1}\cdots\psi_{i_q}$ for  $I=\{i_1<\cdots<i_q\}$. Hence, the above identity yields
\begin{align*}
  U_s\psi_IU_s^*
  &=
  \mathrm i^{q/2}
  (U_s\psi_{i_1}U_s^*)\cdots
  (U_s\psi_{i_q}U_s^*)\\
  &=
  \left(\prod_{x\in I}s_x\right)\psi_I
  =\chi_s(I)\psi_I.
\end{align*}
Writing $g^{(s)}:=(\chi_s(I)g_I)_I$, we therefore have
\begin{align*}
  H_N(g^{(s)})
  &=\sigma_N\sum_{I\in\calI_N}\chi_s(I)g_I\psi_I\\
  &=\sigma_N\sum_{I\in\calI_N}g_IU_s\psi_IU_s^*
    =U_sH_N(G)U_s^*.
\end{align*}
Using the unitary equivariance of the matrix exponential, we have 
\[
  e^{-\beta H_N(g^{(s)})}
  =U_se^{-\beta H_N(G)}U_s^*.
\]
Finally, since the trace is invariant under unitary conjugation, we conclude 
\[
  Z_N(\beta;g^{(s)})
  =\tr\!\left(U_se^{-\beta H_N(G)}U_s^*\right)
  =\tr e^{-\beta H_N(G)}
  =Z_N(\beta;G),
\]
which proves \eqref{eq:sign-gauge-invariance}.
}

Assume that $\alpha$ is odd so that  we can take $x$ such that
$\sum_{I:x\in I}\alpha_I$ is odd.  We take $s$ such that $s_y=1$ for
$y\neq x$ and $s_x=-1$.  Then,
 \[
  H_\alpha((\chi_s(I)g_I)_I)
  =
  \left(
    \prod_{x=1}^N
    s_x^{\textcolor{red}{\sum_{I:x\in I}}\textcolor{blue}{\sum_{I\ni x}}\alpha_I}
  \right)H_\alpha(G) = -H_\alpha(G).
\]
\textcolor{red}{
By~\eqref{eq:signed-permutation-invariance} with $\pi=\mathrm{id}$,
$Z_N(\beta;T_{\mathrm{id},s}G)=Z_N(\beta;G)$.
Since $T_{\mathrm{id},s}G$ and $G$ have the same distribution} \textcolor{blue}{Hence}, we have
\begin{align*}\textcolor{red}{
\E[Z_N(G) H_\alpha(G)] = \E[Z_N(\chi_s G) H_\alpha(\chi_s G)] = 
\E[Z_N(G) H_\alpha(G)]}\\
\textcolor{blue}{
\E[Z_N(G') H_\alpha(G')] =\E[Z_N(\chi_s G') H_\alpha(\chi_s G')]=\E[Z_N(G') H_\alpha(G')]}, 
\end{align*}
which gives the second claim. 
\end{proof}

\textcolor{blue}{\subsection{Gaussian integration by parts}}

For a multi-index $\alpha \in \{0,1,\ldots\}^{\mathcal I_N}$, we denote the associated multi-index partial derivative of a function $f\colon \mathbb R^{\mathcal I_N} \to \mathbb R$ as $\partial^\alpha f$.

\begin{lemma}\label{lem:Gaussian-Hermite-IBP}
Let $f\colon\mathbb R^{\mathcal I_N}\to\mathbb R$ be of class $C^{|\alpha|}$,
and suppose that $f$ and its multi-index partial derivatives up to order $\alpha$
have at most exponential growth in \textcolor{red}{$|G|$} \textcolor{blue}{$\norm{G}$}\textcolor{red}{, where $|G|$ denotes the Euclidean norm}: for any multi-index $\eta=(\eta_I)_{I\in\mathcal I_N}$ with $\eta_I\leq \alpha_I$, there exists a  constant $C_\eta>0$ such that
\begin{equation*}
  |(\partial^\eta f)(G)| \leq \textcolor{red}{C_\eta e^{C_\eta |G|}}\ \textcolor{blue}{C_\eta e^{C_\eta \|G\|}}.
\end{equation*}
Then
\begin{equation}\label{eq:Gaussian-Hermite-IBP}
  \E\!\left[H_\alpha(G)f(G)\right]
  =\E\!\left[\partial_g^\alpha f(G)\right].
\end{equation}
In particular, \eqref{eq:Gaussian-Hermite-IBP} holds for
$f(G)=Z_N(\beta;G)$.
\end{lemma}
\begin{proof}
Recall that the one-dimensional integration-by-parts identity yields 
\[\textcolor{red}{
\E\!\left[\operatorname{He}_k(g)h(g)\right] = \E\!\left[h^{(k)}(g)\right]}\ \textcolor{blue}{
\E\!\left[\operatorname{He}_k(G)f(G)\right] = \E\!\left[f^{(k)}(G)\right]},
\]
for every function $\textcolor{red}{h}\ \textcolor{blue}{f}\in C^k(\mathbb R)$ satisfying the stated growth
condition~\cite[Equation~(2.3.15)]{S2023}. \textcolor{red}{Here, $g$ is a standard Gaussian random variable.} \textcolor{red}{By the definition of $H_\alpha(G)$, applying} \textcolor{blue}{Applying} this identity successively to the independent Gaussian
coordinates $(g_I)_{I\in\mathcal I_N}$ \textcolor{blue}{and using
\[
H_\alpha(G)
=\prod_{I\in\mathcal I_N}\operatorname{He}_{\alpha_I}(g_I),
\]}
gives
\begin{align*}
\E\!\left[
H_\alpha(G)f(G)
\right]
&=
\E\!\left[
\left(
\prod_{I\in\mathcal I_N}
\operatorname{He}_{\alpha_I}(g_I)
\right)
f(G)
\right] \\
&=
\E\!\left[
\left(
\prod_{I\in\mathcal I_N}
\partial_{g_I}^{\alpha_I}
\right)
f(G)
\right] \\
&=
\E\!\left[
\partial_g^\alpha f(G)
\right].
\end{align*}
The growth assumption justifies all integrations by parts. Finally, for fixed
$N$ and $\beta$, the partition function $Z_N(\beta;G)$ is a smooth function
of finitely many Gaussian coordinates, and all of its derivatives have at
most exponential growth in  \textcolor{red}{$|G|$} \textcolor{blue}{$\norm{G}$}.  
Thus it satisfies the assumption.
\end{proof}
\textcolor{blue}{\subsection{First and second moment}
\begin{lemma}\label{lem:exact-second-chaos}
For every $\beta\in \mathbb R$,
\[
  P_1Y_N(\beta)=0
\]
and
\begin{equation}\label{eq:P2Y-exact}
  P_2 Y_N(\beta)
  =
\frac{\beta\,\partial_\beta\log \E Z_N(\beta)}{2M_N}
  \sum_{I\in\calI_N}(g_I^2-1).
\end{equation}
\end{lemma}}

\begin{proof}[{Proof \textcolor{red}{ of Proposition \ref{prop:exact-second-chaos}}}]
For \textcolor{red}{any} multi-index $\alpha$ with $|\alpha|=1$, since $\alpha$ is always odd, the first claim follows from Lemma~\ref{lem: parity rule lemma}. 

\textcolor{red}{Note that, if a multi-index $\alpha$ satisfies $|\alpha|=2$ and $\alpha_I=\alpha_J=1$ for some $I,J\in\calI_N$ with $I\ne J$, then $\sum_{K: x\in K} \alpha_K=1$ for any $x\in I\setminus J$, and hence $\alpha$ is odd.
Fix any $J,J'\in\calI_N$ and choose $\pi\in\mathfrak S_N$
such that $\pi(J)=J'$.
Let $G_\pi:=T_{\pi,\mathbf 1}G$, where $\mathbf 1=(1,\ldots,1)$.
By~\eqref{eq:signed-permutation-invariance} with $s=\mathbf 1$,
$Z_N(\beta;G_\pi)=Z_N(\beta;G)$.
Since $G$ and $G_\pi$ have the same distribution, we obtain
\begin{align}\label{eq: second-coeffiicient-invariance}
  \E[Z_N(\beta;G)(g_J^2-1)]
  &=\E\!\left[Z_N(\beta;G_\pi)
    \bigl((\epsilon_\pi(J)g_{\pi(J)})^2-1\bigr)\right]\notag\\
  &=\E[Z_N(\beta;G)(g_{J'}^2-1)].
\end{align}}
\textcolor{blue}{
  In the second chaos it also kills every 
monomial $g_Ig_J$ with $I\neq J$, since the corresponding multi-index is also odd.  Consequently, since ${\rm He}_2(x)=x^2-1$, 
\[
  P_2Z_N(\beta)=\frac{1}{2} \sum_I \E[Z_N(\beta;G') ((g'_I)^2-1) ] (g_I^2-1).
\]
The expectation
\(
\E[Z_N(\beta;G)(g_I^2-1)]
\)
is independent of $I$.  Indeed,
a permutation of the sites sends any
$q$-set to any other one; the induced signs on the Majorana monomials are
absorbed by the coordinatewise sign symmetry of the independent centered
Gaussian couplings, while the normalized trace is invariant under the
corresponding Clifford automorphism. It therefore remains to compute the
common coefficient
\[
b_N(\beta):= \frac{1}{2}\E[Z_N(\beta;G) (g_I^2-1)].
\]
Applying Gaussian integration by parts yields
\[
\E\bigl[Z_N(\beta)(g_I^2-1)\bigr]
=\E\,g_I \partial_{g_I} Z_N(\beta).
\]
}

On the other hand,  since $
  Z_N(\beta;tG)
  =Z_N(t\beta;G),
$ 
 differentiating at $t=1$ gives
 \begin{equation}\label{eq:partial-derivative-relaton}
 \sum_Ig_I\partial_{g_I}Z_N(\beta;G)
=\beta\,\partial_\beta Z_N(\beta;G).
 \end{equation}
\textcolor{red}{Therefore, substituting ${\rm He}_2(x)=x^2-1$ we obtain
\begin{align*}
	P_2Z_N(\beta)&\overset{\eqref{eq:odd-coefficient-zero}}{=}\sum_{I\in\calI_N}\frac{\E[Z_N(\beta;G)(g_I^2-1)]}{2}(g_I^2-1)\\
	&\overset{\eqref{eq: second-coeffiicient-invariance}}{=}\frac{1}{2|\calI_N|}\E\left[\sum_{I\in\calI_N}Z_N(\beta;G)(g_I^2-1)\right]\sum_{I\in\calI_N}(g_I^2-1)\\
	&=\frac{1}{2|\calI_N|}\E\left[\sum_{I\in\calI_N}g_I\partial_{g_I}Z_N(\beta;G)\right]\sum_{I\in\calI_N}(g_I^2-1)\\
	&\overset{\eqref{eq:partial-derivative-relaton}}{=}\frac{\beta\E[\partial_\beta Z_N(\beta;G)]}{2|\calI_N|}\sum_{I\in\calI_N}(g_I^2-1)\\
	&=\frac{\beta\partial_\beta\E[ Z_N(\beta;G)]}{2|\calI_N|}\sum_{I\in\calI_N}(g_I^2-1)
\end{align*}
where the third equality follows from the (one-dimensional) Gaussian integration by parts, and the last equality follows from the dominated convergence theorem because
\[\sup_{|\beta|<B}|\partial_\beta Z_N(\beta;G)|\le \|H_N(G)\|_\infty e^{B\|H_N(G)\|_\infty},\quad  \|H_N(G)\|_\infty\le \sigma_N\sum_{I\in\calI_N}|g_I|\]
for every $B>0$, and the bound is Gaussian integrable.} \textcolor{blue}{
Hence, we have
\begin{align*}
  2M_Nb_N(\beta)
  & =\sum_{I\in\calI_N}\E\bigl[g_I\partial_{g_I}Z_N(\beta)\bigr]=\beta\,\partial_\beta\E Z_N(\beta).
\end{align*}}
This proves \textcolor{red}{$P_2 Y_N(\beta)
	=A_N(\beta)Q_N$, because $P_2Y_N(\beta)=P_2Z_N(\beta)/\E Z_N(\beta)$} \textcolor{blue}{\eqref{eq:P2Y-exact}}.
\end{proof}

\end{document}
